% ЕМТ 2022, XII клас — проверен LaTeX препис на липсващите задачи 12.1 и 12.3 от официалните файлове в Klasirane.com.
% Условия: https://klasirane.com/api/competitions/EMT/All/2022/12%20%D0%BA%D0%BB/probs
% SHA-256 (условия): f2736e5a6043d0f81ac8e1a370f72f9d60c2793a989cacafc5d16c39328ae2fb
% Решения: https://klasirane.com/api/competitions/EMT/All/2022/12%20%D0%BA%D0%BB/sol
% SHA-256 (решения): 7ce18f9de96669768e1dac52411a263261ab9c1b3c25f469fbe5bd825aa2eb70
% Mathpix Snip (условия): 728b37f0-fb0e-4297-99d2-a17ac77c7dcd
% Mathpix Snip (решения): 3dca50ae-1167-4546-adb4-c3658e085efa
% Формулите, формулировките и точковите оценки са сверени визуално с официалните PDF страници.
% Картинният маркер на Mathpix между 12.3 и решението е чертежът към 12.2 от началото на следващата PDF страница; 12.1 и 12.3 нямат фигури.

Задача 12.1. Нека $x+y+z=1$, $x^{2}+y^{2}+z^{2}=2$ и $x^{3}+y^{3}+z^{3}=3$. Да се намери стойността на израза $A=x^{5}+y^{5}+z^{5}$.

Решение. Нека $x, y, z$ са корени на полинома от трета степен $P(t)=t^{3}+a t^{2}+b t+c=(t-x)(t-y)(t-z)$. Тогава по формулите на Виет имаме $a=-(x+y+z)=-1$, $b=x y+y z+z x$, $c=-x y z$.

От равенството $(x+y+z)^{2}=x^{2}+y^{2}+z^{2}+2(x y+y z+z x)$ получаваме $1^{2}=2+2 b$, т.е. $b=-\frac{1}{2}$.

Нека означим с $S_{k}=x^{k}+y^{k}+z^{k}$ степенните сборове. Имаме $S_{0}=3$, $S_{1}=1$, $S_{2}=2$, $S_{3}=3$. Сега от формулата на Нютон $S_{k+3}+a S_{k+2}+b S_{k+1}+c S_{k}=0$ (формулата се получава след умножаване на всяко от равенствата $x^{3}+a x^{2}+b x+c=0$, $y^{3}+a y^{2}+b y+c=0$, $z^{3}+a z^{2}+b z+c=0$ съответно с $x^{k}$, $y^{k}$, $z^{k}$ и почленното им събиране) при $k=0$ имаме $3+(-1) \cdot 2+\left(-\frac{1}{2}\right) \cdot 1+c \cdot 3=0$, т.е. $c=-\frac{1}{6}$. Така последователно получаваме
$$
S_{4}=S_{3}+\frac{1}{2} S_{2}+\frac{1}{6} S_{1}=3+\frac{1}{2} \cdot 2+\frac{1}{6} \cdot 1=\frac{25}{6},
$$
$$
S_{5}=S_{4}+\frac{1}{2} S_{3}+\frac{1}{6} S_{2}=\frac{25}{6}+\frac{1}{2} \cdot 3+\frac{1}{6} \cdot 2=\frac{25+9+2}{6}=6.
$$

Оценяване. (6 точки) 2 т. за свеждане до полином от трета степен с нули $x, y, z$; 2 т. за намиране на коефициентите на полинома; 2 т. за намиране на $A$.

Задача 12.3. Редицата $a_{n}$ е зададена чрез $a_{1} \geq 2$ и рекурентната връзка
$$
a_{n+1}=a_{n} \sqrt{\frac{a_{n}^{3}+2}{2\left(a_{n}^{3}+1\right)}}
$$
за $n \geq 1$. Да се докаже, че за всяко естествено число $n$ е в сила неравенството $a_{n}>\sqrt{\frac{3}{n}}$.

Решение. Да забележим, че равенството е еквивалентно на
$$
\frac{1}{a_{n+1}^{2}}-\frac{1}{a_{n}^{2}}=\frac{a_{n}}{a_{n}^{3}+2}. \tag{*}
$$
От СА-СГ следва, че $\frac{a_{n}}{a_{n}^{3}+2} \leq \frac{1}{3}$. Така получаваме $\frac{1}{a_{n+1}^{2}}-\frac{1}{a_{n}^{2}} \leq \frac{1}{3}$ за всяко $n \in \mathbb{N}$, което след сумиране дава $\frac{1}{a_{n}^{2}} \leq \frac{n-1}{3}+\frac{1}{a_{1}^{2}}<\frac{n}{3}$, откъдето получаваме желания резултат.

Оценяване. (7 точки) 3 т. за (*), 1 т. за $\frac{a_{n}}{a_{n}^{3}+2} \leq \frac{1}{3}$, 1 т. за $\frac{1}{a_{n+1}^{2}}-\frac{1}{a_{n}^{2}} \leq \frac{1}{3}$ и 2 т. за довършване.
